\documentclass[11pt,a4paper]{article}

\usepackage[utf8]{inputenc}
\usepackage[T1]{fontenc}
\usepackage{amsmath,amsfonts,amssymb}
\usepackage{booktabs}
\usepackage[margin=2.5cm]{geometry}
\usepackage{hyperref}
\usepackage{enumitem}

\title{Case Studies CS1--CS4:\\Stochastic Lorenz-63 Data Assimilation}
\author{4DVarNet-FM}
\date{}

\begin{document}
\maketitle

\section{State-Space Model}

We consider a stochastic Lorenz-63 system augmented with a fourth variable $W_L$ representing a wind forcing. The state vector is
\[
\mathbf{x} = [X, Y, Z, W_L]^{\mathsf T} \in \mathbb{R}^{4}.
\]

The deterministic drift is given by
\begin{align}
\dot X &= \sigma\,(Y - X) + c_1\,W_L,\\
\dot Y &= X\,(\rho - Z) - Y,\\
\dot Z &= X\,Y - \beta\,Z,\\
\dot W_L &= -\gamma\,(W_L - \overline W_L) + c_2\,X,
\end{align}
with parameters $\sigma = 10.0$, $\rho = 28.0$, $\beta = 8/3$, $\gamma = 0.05$, $\overline W_L = 0.0$, $c_1 = 1.0$, $c_2 = 0.1$.

The system is discretised with the Euler--Maruyama scheme using a time step $\Delta t = 0.01$. Multiplicative noise is injected into $Y$, $Z$ and $W_L$:
\begin{align}
X_{t+1} &= X_t + \dot X_t\,\Delta t,\\
Y_{t+1} &= Y_t + \dot Y_t\,\Delta t + \sigma_0\,Y_t\,\sqrt{\Delta t}\,\xi_1,\\
Z_{t+1} &= Z_t + \dot Z_t\,\Delta t + \sigma_0\,Z_t\,\sqrt{\Delta t}\,\xi_2,\\
W_{L,t+1} &= W_{L,t} + \dot W_{L,t}\,\Delta t + \sigma_L\,\sqrt{\Delta t}\,\xi_3,
\end{align}
where $\xi_1, \xi_2, \xi_3 \sim \mathcal{N}(0,1)$ are independent Gaussian increments, $\sigma_0 = 0.08$ and $\sigma_L = 0.20$.

\section{Observation Model}

Sparse noisy observations of the state variables $(X,Y,Z)$ are available every $\Delta t_{\text{obs}}=20\Delta t = 0.2$ time units:
\[
\mathbf{y}_k = H\,\mathbf{x}(t_k) + \boldsymbol\varepsilon_k,\qquad
\boldsymbol\varepsilon_k \sim \mathcal{N}\bigl(\mathbf{0}, R\bigr),
\]
with $R = R_{\text{var}}\,I_3$ and $R_{\text{var}} = 0.5$. The observation operator $H\in\mathbb{R}^{3\times4}$ selects the first three components. The observation mask is
\[
m(t) = \begin{cases}
1 & \text{if } t \in \{\Delta t_{\text{obs}},\,2\Delta t_{\text{oss}},\,\dots,\,T_{\max}\},\\
0 & \text{otherwise}.
\end{cases}
\]

\section{Case Study CS1: Perfect Model Baseline}

CS1 uses the true dynamics with correct parameters and uncorrupted forcing:
\begin{itemize}[noitemsep]
\item Parameters: $\sigma = 10.0$, $\rho = 28.0$, $\beta = 8/3$ (no bias).
\item Forcing: $W_L^* = W_L$ (the true wind, uncorrupted).
\item Coupling in $\dot X$: linear term $c_1 W_L$.
\item Expected: good reconstruction with RMSE $\ll 0.5$.
\end{itemize}

\section{Case Study CS2: Model Mismatch}

CS2 introduces three sources of model mismatch.

\subsection{Parameter Bias}

The Lorenz parameters are biased by a fraction $b = 0.15$:
\begin{align}
\sigma_{\text{CS2}} &= \sigma\,(1 - b) = 8.5,\\
\rho_{\text{CS2}}  &= \rho\,(1 - b) = 23.8,\\
\beta_{\text{CS2}} &= \beta\,(1 + b) \approx 3.11.
\end{align}

\subsection{Forcing Corruption}

The true wind $W_L$ is corrupted by an Ornstein--Uhlenbeck process $\eta$ and a state-dependent term:
\[
W_L^* = W_L + \eta + \alpha\,X,
\]
where $\alpha = 0.15$ (forcing state bias). The corruption $\eta$ satisfies
\[
d\eta = -\frac{1}{\tau_\eta}\,\eta\,dt + \sigma_\eta\sqrt{\frac{2}{\tau_\eta}}\,dW,
\]
with $\tau_\eta = 5.0$ and $\sigma_\eta = \sqrt{0.5}$.

\subsection{Nonlinear Coupling}

The coupling in $\dot X$ is changed from linear to quartic:
\[
c_1\,W_L \;\longrightarrow\; \operatorname{sgn}(W_L)\,W_L^2,
\]
which introduces a stronger and non-symmetric forcing on the $X$ variable.

\section{Case Study CS3: Randomised Parameters (Linear Coupling)}

CS3 applies the same parameter randomisation used during $E3/F3$ training to the CS1 configuration (linear coupling, no forcing corruption). Each test window draws random parameters:

\[
\sigma \sim \mathcal{U}\bigl(\sigma_{\text{true}}(1 - \epsilon),\;\sigma_{\text{true}}(1 + \epsilon)\bigr),\quad
\rho \sim \mathcal{U}\bigl(\rho_{\text{true}}(1 - \epsilon),\;\rho_{\text{true}}(1 + \epsilon)\bigr),\quad
\beta \sim \mathcal{U}\bigl(\beta_{\text{true}}(1 - \epsilon),\;\beta_{\text{true}}(1 + \epsilon)\bigr),
\]

with $\epsilon = 0.2$ (20\% noise), independently per window. The forcing is the true wind $W_L$ (uncorrupted) and the coupling in $\dot X$ remains linear.

\section{Case Study CS4: Randomised Parameters (Quartic Coupling)}

CS4 combines the CS2 model mismatch (parameter bias, OU-corrupted forcing, quartic coupling) with the same per-window parameter randomisation used in CS3:

\begin{itemize}[noitemsep]
\item Parameter bias $b = 0.15$ applied after randomisation.
\item Forcing corrupted by OU process $\eta$ and state-dependent term $\alpha X$ ($\alpha = 0.15$).
\item Coupling: quartic $\operatorname{sgn}(W_L)\,W_L^2$.
\item Per-window parameter noise $\epsilon = 0.2$ (same as CS3).
\end{itemize}

This is the most challenging case: the model must handle simultaneously (i) biased parameters, (ii) corrupted forcing, (iii) nonlinear coupling, and (iv) unseen random parameter draws.

\section{Parameter Summary}

\begin{table}[ht]
\centering
\caption{Default model parameters.}
\begin{tabular}{lccr}
\toprule
Symbol & Value & Unit & Description\\
\midrule
$\sigma$ & $10.0$ & — & Lorenz $X\leftrightarrow Y$ feedback\\
$\rho$ & $28.0$ & — & Rayleigh number\\
$\beta$ & $8/3$ & — & Aspect ratio\\
$\gamma$ & $0.05$ & s$^{-1}$ & Wind relaxation rate\\
$\overline W_L$ & $0.0$ & — & Wind baseline\\
$c_1$ & $1.0$ & — & Wind $\to$ state coupling\\
$c_2$ & $0.1$ & — & State $\to$ wind coupling\\
$\sigma_0$ & $0.08$ & — & Multiplicative noise std\\
$\sigma_L$ & $0.20$ & — & Wind noise std\\
$\tau_\eta$ & $5.0$ & s & OU correlation time\\
$\sigma_\eta$ & $\sqrt{0.5}$ & — & OU noise std\\
$\alpha$ & $0.15$ & — & Forcing-state bias\\
$b$ & $0.15$ & — & Parameter bias fraction\\
$R_{\text{var}}$ & $0.5$ & — & Observation noise variance\\
$B_{\text{var}}$ & $2.0$ & — & Background variance\\
$\Delta t$ & $0.01$ & s & Integration step\\
$T_{\max}$ & $3.0$ & s & Assimilation window\\
$\Delta t_{\text{obs}}$ & $0.2$ & s & Observation interval\\
\bottomrule
\end{tabular}
\end{table}

\end{document}
